\section{Abschlussinterview nach den Usability-Tests}
\label{sec:BB-06_feedbackinterview}

\section{Kodierte Rohdaten der leitfadengestuetzten Interviews}
\label{sec:BB-05_interviews_kodiert}

Die folgenden Tabellen und Kategorien fassen die Interviewaussagen aller zehn Teilnehmenden zusammen. 
Die Struktur folgt dem originalen Leitfaden (Bloecke A--D). 
Qualitative Aussagen sind thematisch gebuendelt und mit den zugehoerigen TN-IDs gekennzeichnet.

% =====================================================
\subsection{Block A -- Wahrgenommene Nuetzlichkeit und Gebrauchstauglichkeit (RQ-1)}

\subsubsection{GF-1 -- Gesamtnutzen neues Wiki (Skala)}

\begin{table}[H]
\centering
\caption{Bewertungen des Gesamtnutzens (GF-1)}
\label{tab:BB-gf1}
\begin{tabular}{c c}
\toprule
TN-ID & Bewertung (1--5) \\
\midrule
TN-01 & 5 \\
TN-02 & 4 \\
TN-03 & 4 \\
TN-04 & 4 \\
TN-05 & 4 \\
TN-06 & 5 \\
TN-07 & 5 \\
TN-08 & 4 \\
TN-09 & 5 \\
TN-10 & 5 \\
\bottomrule
\end{tabular}
\end{table}

\subsubsection{HF-1 -- Begruendung Gesamtnutzen}

\textbf{Verbesserte Auffindbarkeit und Effizienz}
\begin{itemize}
  \item Inhalte schneller auffindbar, weniger Klicks notwendig (TN-01, TN-02, TN-04, TN-06, TN-07, TN-09).
  \item Klarere Navigation und logischere Menuestruktur (TN-01, TN-03, TN-05, TN-10).
\end{itemize}

\textbf{Visuelle Klarheit und Reduktion}
\begin{itemize}
  \item Reduziertes, modernes Design erleichtert Orientierung (TN-01, TN-07, TN-09).
\end{itemize}

\subsubsection{GF-2 -- Verstaendlichkeit und Orientierung (Skala)}

\begin{table}[H]
\centering
\caption{Bewertungen der Orientierung (GF-2)}
\label{tab:BB-gf2}
\begin{tabular}{c c}
\toprule
TN-ID & Bewertung (1--5) \\
\midrule
TN-01 & 4 \\
TN-02 & 4 \\
TN-03 & 4 \\
TN-04 & 4 \\
TN-05 & 4 \\
TN-06 & 5 \\
TN-07 & 5 \\
TN-08 & 4 \\
TN-09 & 5 \\
TN-10 & 5 \\
\bottomrule
\end{tabular}
\end{table}

\subsubsection{HF-2 -- Mentale Modelle und Informationsarchitektur}

\textbf{Uebereinstimmung mit mentalen Modellen}
\begin{itemize}
  \item Inhalte dort auffindbar, wo sie intuitiv erwartet werden (TN-01, TN-02, TN-03, TN-06, TN-08, TN-09, TN-10).
  \item Klare thematische Buendelung und nachvollziehbare Struktur (TN-01, TN-04, TN-07).
\end{itemize}

\textbf{Bruchstellen}
\begin{itemize}
  \item Ansprechpartner schwer auffindbar oder inkonsistent platziert (TN-01, TN-02, TN-05, TN-06, TN-07, TN-09, TN-10).
  \item Nachrichten auf der Startseite zu dominant (TN-03, TN-04, TN-05, TN-07).
  \item Unklare Scroll-Logik auf langen Seiten (TN-03, TN-06, TN-08).
\end{itemize}

% =====================================================
\subsection{Block B -- Personas, Rollenpassung und UX-Methoden (RQ-2)}

\subsubsection{GF-3 -- Passung zur Rolle / Persona (Skala)}

\begin{table}[H]
\centering
\caption{Bewertungen der Rollenpassung (GF-3)}
\label{tab:BB-gf3}
\begin{tabular}{c c}
\toprule
TN-ID & Bewertung (1--5) \\
\midrule
TN-01 & 5 \\
TN-02 & 4 \\
TN-03 & 4 \\
TN-04 & 4 \\
TN-05 & 4 \\
TN-06 & 5 \\
TN-07 & 4 \\
TN-08 & 4 \\
TN-09 & 5 \\
TN-10 & 5 \\
\bottomrule
\end{tabular}
\end{table}

\subsubsection{HF-3 -- Konkrete Strukturpassung}

\textbf{Positive Aspekte}
\begin{itemize}
  \item Alle relevanten Inhalte fuer die jeweilige Rolle vorhanden (TN-01, TN-02, TN-04, TN-05, TN-06, TN-09, TN-10).
  \item Einstiegs- und Uebersichtsseiten erleichtern Orientierung (TN-01, TN-02, TN-04, TN-06).
\end{itemize}

\textbf{Optimierungswuensche}
\begin{itemize}
  \item Wunsch nach staerkerer Sichtbarkeit rollenrelevanter Inhalte (TN-03, TN-07, TN-10).
  \item Differenzierung nach Erfahrungsstufen oder Rollen (TN-03, TN-07).
\end{itemize}

\subsubsection{OF-1 -- Wahrgenommene UX-Methoden}

\begin{itemize}
  \item Nutzerzentrierte Gestaltung erkennbar durch klare Struktur und konsistente Benennung (TN-01, TN-02, TN-06, TN-09).
  \item Generische Struktur fuer alle Rollen akzeptiert (TN-04, TN-05, TN-08).
\end{itemize}

% =====================================================
\subsection{Block C -- Akzeptanz, Nutzungsintention und Kommunikation (RQ-3)}

\subsubsection{GF-4 -- Nutzungsintention (Skala)}

\begin{table}[H]
\centering
\caption{Bewertungen der Nutzungsintention (GF-4)}
\label{tab:BB-gf4}
\begin{tabular}{c c}
\toprule
TN-ID & Bewertung (1--5) \\
\midrule
TN-01 & 5 \\
TN-02 & 4 \\
TN-03 & 4 \\
TN-04 & 4 \\
TN-05 & 4 \\
TN-06 & 5 \\
TN-07 & 4 \\
TN-08 & 3 \\
TN-09 & 4 \\
TN-10 & 4 \\
\bottomrule
\end{tabular}
\end{table}

\subsubsection{HF-4 -- Akzeptanztreiber und -barrieren}

\textbf{Akzeptanztreiber}
\begin{itemize}
  \item Aktualitaet und Verlaesslichkeit der Inhalte (TN-01, TN-02, TN-04, TN-05, TN-09, TN-10).
  \item Zentrale Wissensquelle statt paralleler Ablagen (TN-01, TN-02, TN-06, TN-07, TN-10).
  \item Gute Struktur und Auffindbarkeit (TN-01, TN-03, TN-06, TN-08, TN-09).
  \item Konkreter Nutzen im Arbeitsalltag (TN-01, TN-04, TN-06, TN-07).
\end{itemize}

\textbf{Akzeptanzbarrieren}
\begin{itemize}
  \item Veraltete oder unvollstaendige Inhalte (TN-02, TN-04, TN-05, TN-07, TN-09, TN-10).
  \item Schwache Suchfunktion (TN-01, TN-03, TN-06, TN-07, TN-09).
  \item Mehrere parallele Wissensquellen (TN-01, TN-02, TN-05, TN-08).
\end{itemize}

\subsubsection{HF-5 -- Kommunikations- und Einfuehrungserwartungen}

\begin{itemize}
  \item Vorstellung in bestehenden Meetings (TN-01, TN-02, TN-04, TN-06, TN-07, TN-08, TN-09, TN-10).
  \item Dialogische Austauschformate und Q\&A (TN-01, TN-03, TN-06, TN-07, TN-09, TN-10).
  \item Regelmaessige Updates oder Newsletter (TN-01, TN-04, TN-05, TN-07, TN-08, TN-09).
  \item Rolle von Fuehrungskraeften als Multiplikatoren (TN-02, TN-06, TN-07, TN-10).
\end{itemize}

\subsubsection{OF-2 -- Beitrag zur Zusammenarbeit}

\begin{itemize}
  \item Verbesserung von Transparenz und gemeinsamer Wissensbasis (TN-01, TN-02, TN-06, TN-07, TN-10).
  \item Grenzen durch inkonsistente Pflege oder konkurrierende Systeme (TN-01, TN-03, TN-05, TN-08).
\end{itemize}

% =====================================================
\subsection{Block D -- Offenes Feedback}

\subsubsection{OF-3 -- Offene Eindruecke}

\begin{itemize}
  \item Positiv: Modernes Design, uebersichtliche Startseite, klare visuelle Struktur (TN-01, TN-07, TN-09).
  \item Keine oder kaum negative Ueberraschungen genannt (Mehrheit).
\end{itemize}

\subsubsection{OF-4 -- Wichtigste Verbesserung}

\begin{itemize}
  \item Wunsch nach Chatbot oder intelligenter Suchunterstuetzung (TN-01, TN-05, TN-07, TN-09).
  \item Verbesserung der Startseiten-Priorisierung (TN-03, TN-04, TN-05, TN-07).
  \item Ausbau der Suchfunktion (TN-02, TN-06, TN-08).
\end{itemize}
